\scanpage{43}
D'autre part, on a
\[
R(D/U\ k)\hookrightarrow R(D,k)
\]
et l'image de $R(D/U,k)$ dans $R(D,k)$ est précisément formée des classes \unsure{des|de} $E$ satisfaisant la condition du début. Pour un tel $E$, on a donc une opération $\unsure{\natural|\xi}$ sur la \uline{classe} de $E$. Donc $\operatorname{cl}(E)^{\unsure{\natural|\xi}}$ est défini dès qu'il existe un ss-grp \deleted{\unsure{distingué|d'indice}} $U$ d'indice fini dans $I$ tel que \deleted{$D/U$} $U$ opère sur $E$ de façon unipotente. / Si on prend des classes de repr. \uline{continues} ($k$ étant un corps top) on \deleted{prend} \inserted{obtient} des \deleted{$\unsure{=}$} $\operatorname{cl}(E)^{\unsure{\natural|\xi}}$ comme classes de repr. \uline{continues} de $D/I$. Compatibilité avec ext des corps de base.

\deleted{Supposons que $D/I\simeq\mathbf{Z}$, et \unsure{donc} qu'une classe de repr. sur $D/I$ est donnée par sa série car. - $L_{E^{\unsure{\natural|\xi}}}(t)$. On aura alors}

\margin{42}

